\section{实验与结果分析}

本节检验第5节方法的有效性，围绕四个问题展开：其一，在中心场景上，本文方法相对经典规则、库内最优单规则、遗传规划生成的规则、滚动重排程搜索方法、其他学习型规则选择方法以及直接派工的学习方法表现如何；其二，增益来自规则库、理论特征还是网络结构，学到的规则选择是否与状态相关；其三，在中心场景训练的策略，换到机器数、负载与交期松紧不同的场景，或改变检查点数量、预测出现偏差时，能否保持优势；其四，收敛过程与计算开销如何。6.1 节给出实验设置，6.2 节说明学习率与折扣因子的预实验，6.3 节为中心场景主对比，6.4 节为消融实验与规则选择行为分析，6.5 节为未见场景上的泛化与压力测试，6.6 节报告收敛与运行时间，6.7 节给出公开真实工业订单轨迹驱动的案例研究。

%=====================================================================
\subsection{实验设置}

\subsubsection{实例场景}

实例按第3节的问题结构随机生成。车间有 $M$ 台机器，$t=0$ 时已有 $N_0=20$ 个作业，此后共有 $N_{\mathrm{new}}$ 个新作业依次到达，到达间隔独立服从均值为 $1/\lambda$ 的指数分布；第 $N_{\mathrm{new}}$ 个新作业到达后不再产生到达，即到达过程是达到规定次数后停止的泊松过程。每个作业的工序数在 2 至 5 之间均匀取整数，每道工序的合格机器数在 1 至 $M$ 之间均匀取整数，加工时间在 $[1,50]$ 上均匀取值，作业的服务数量 $q_j$ 独立取自离散均匀分布 $U\{1,2,3,4,5\}$，因此 $\mathbb{E}[q]=3$、$\mathbb{E}[q^2]=11$。作业的交货期为到达时刻加上交期紧度与其各工序平均加工时间之和的乘积。平均到达间隔 $1/\lambda$ 由目标利用率反推，即单个作业的平均加工总量除以目标利用率与机器数的乘积。以上设定沿用动态柔性作业车间调度常用的基准生成方式~\cite{luo2020,tiacci2024}；加工时间的下界取 1，以保证加工时间为正。特征归一化尺度取自预先公开的生成分布：$\bar p=\mathbb{E}[p]=25.5$ 为单道工序的名义平均加工时间，$\bar W=\mathbb{E}[L]\bar p=3.5\times25.5=89.25$ 为单个作业的名义平均总加工量。二者不按某一算例实际抽样的作业均值重算，也不等同于各工序在合格机器上的最短加工时间之和；在全部场景中固定。记目标利用率为 $\upsilon$，则到达率设为 $\lambda=\upsilon M/\bar W$，仅由公开的场景参数确定。

场景由四个因子确定（表~\ref{tab:factors}）：机器数、目标利用率、作业交期紧度和服务权重倍率。机器数的水平参照文献中常用的 5、10、20 台~\cite{li2025,zhang2024}；交期紧度取 0.5、1.0、1.5，与 Luo 基准及其后续工作一致~\cite{tiacci2024,lei2024}；目标利用率取 0.65、0.80、0.95，取自 DynaSchedBench 的难度网格~\cite{dynasched2026}。全部场景固定 $N_0=20$、$N_{\mathrm{new}}=50$，共70个作业。服务权重倍率 $s_w\in\{0.2,1,5\}$ 是本研究自行设置的目标权衡敏感性因子，而非文献中的作业规模因子；其作用仅是将两个服务检查点的校准权重统一乘以 $s_w$，不改变 $TT$ 的全局权重、服务承诺、检查点时限或实际供给。各算法只在相同场景与相同目标权重下配对比较，不把不同 $s_w$ 的绝对 $F$ 值直接作为跨场景优劣依据。

\begin{table}[t]
  \centering
  \caption{场景因子与水平}\label{tab:factors}
  \small
  \begin{tabular}{@{}>{\raggedright}p{3.0cm}>{\centering}p{2.8cm}>{\raggedright\arraybackslash}p{8.2cm}@{}}
    \toprule
    因子 & 水平 & 参考 \\
    \midrule
    机器数 $M$ & 5、10、20 & Li 等~\cite{li2025}；Zhang 等~\cite{zhang2024} \\
    服务权重倍率 $s_w$ & 0.2、1、5 & 本文目标权衡敏感性设计 \\
    目标利用率 & 0.65、0.80、0.95 & DynaSchedBench~\cite{dynasched2026} \\
    交期紧度 & 0.5、1.0、1.5 & Luo 基准~\cite{tiacci2024}；Lei 等~\cite{lei2024} \\
    \bottomrule
  \end{tabular}
\end{table}

服务部分是本文问题特有的设定，文献中没有直接对应的基准。服务实体数取 $R=|\mathcal{R}^{\mathrm{svc}}|=\max\{1,\operatorname{round}[(N_0+N_{\mathrm{new}})/8]\}$。在实际作业属性和到达轨迹抽样之前，使用独立随机源为每个实体生成正的需求规模权重 $a_r$，再归一化为公开的预期需求份额
\begin{equation}\label{eq:entity_share}
  \pi_r=\frac{a_r}{\sum_{s\in\mathcal{R}^{\mathrm{svc}}}a_s},\qquad \pi_r>0,\quad \sum_r\pi_r=1.
\end{equation}
其中 $a_r\overset{\mathrm{i.i.d.}}{\sim}U[0.8,1.2]$，且每个算例只抽取一次并在整个回合保持不变。每个作业（含初始作业与新到作业）独立按 $\Pr(e_j=r)=\pi_r$ 确定所属实体，不强制每个实体至少分到一个实际作业。公开份额、服务承诺与时限均不依据本实例未来实际释放的作业数量、释放时间、路线、交期或服务数量生成。基准设定每实体有两个累计检查点（$K_r=2$），最终检查点为服务截止时间。
根据公开场景规模与预期份额，实体 $r$ 的期望服务数量及最终承诺定义为
\begin{equation}\label{eq:entity_commitment_gen}
  \bar Q_r=(N_0+N_{\mathrm{new}})\,\mathbb{E}[q] \,\pi_r,\qquad Q_r=\rho^Q\bar Q_r.
\end{equation}
统一的首次累计需求比例固定为 $\eta=0.55$，故 $Q_{r1}=0.55Q_r$，$Q_{r2}=Q_r$；两个检查点的短缺权重固定采用比例 $w_{r1}:w_{r2}=1:3$，即 $w_{r1}=\omega$、$w_{r2}=3\omega$。需求份额 $\pi_r$ 不同，使各实体具有不同的期望服务数量和承诺目标，尽管它们共享同一比例 $\rho^Q$ 和 $\eta$。

检查点时刻在实际作业到达与属性抽样前独立生成。设由公开场景参数决定的基准服务窗口为
\begin{equation}\label{eq:checkpoint_window_gen}
  \bar D=\frac{N_{\mathrm{new}}}{\lambda}+\bar W,
\end{equation}
其中 $N_{\mathrm{new}}/\lambda$ 是计划新到作业到齐时刻的期望跨度，$\bar W$ 是预先固定的名义单作业加工工作量；$\bar D$ 是用于构造承诺时限的日历时间尺度，不代表某个实际排程的完成时刻。每个实体独立抽取 $\xi_r\overset{\mathrm{i.i.d.}}{\sim}U[-0.10,0.10]$，定义
\begin{equation}\label{eq:checkpoint_deadline_gen}
  D_r=\max\{2,\operatorname{round}[\bar D(1+\xi_r)]\},\qquad
  \tau_{r1}=\max\{1,\min\{D_r-1,\operatorname{round}(\theta D_r)\}\},\qquad
  \tau_{r2}=D_r.
\end{equation}
这里对非负数使用四舍五入规则 $\operatorname{round}(x)=\lfloor x+1/2\rfloor$（恰为半整数时向上取整）；上述边界处理保证 $1\le\tau_{r1}<D_r$。同一实体的两个检查点共用其扰动 $\xi_r$，不同实体之间独立；即使一个实体最终没有实际作业，其外生承诺也保持不变。为确保配对试验可复现，实体份额 $a_r$、时限扰动 $\xi_r$、作业属性、到达过程分别使用由实例种子确定的独立随机数流；改变校准参数或预测失配系数时，不重新抽取这些基础随机量。参数 $\rho^Q$、$\theta$ 与统一权重倍率 $\omega$ 在独立校准实例上确定，而 $\eta=0.55$、权重比例 $1:3$ 及扰动范围 $\pm10\%$ 不参与校准。全部检查点的时刻、累计数量与权重在 $t=0$ 公布，并允许实际总供给高于或低于承诺；由此产生的不可避免最终短缺按第3节单独报告。

首次检查点时间比例 $\theta$、承诺比例 $\rho^Q$ 与权重倍率 $\omega$ 仅在独立的30个\emph{中心场景校准实例}（种子25100--25129）上确定。为控制选择自由度，在观察 LLM、GP 或学习型方法的实验结果之前固定候选网格：
\begin{equation}\label{eq:service_calibration_grid}
  \rho^Q\in\{0.70,0.85,1.00,1.15\},\qquad
  \theta\in\{0.45,0.55,0.65\},\qquad
  \omega\in\{0.5,1,2,4,8\}.
\end{equation}
共 $4\times3\times5=60$ 个候选组合。所有组合使用相同的30组作业、实体预期份额及实体时限扰动，仅更改被校准的承诺、首次检查点时刻及权重。其余需求比例、权重比与扰动幅度不参与校准；校准只用 SPT 和 ATC，不涉及 LLM 规则、GP 或学习型方法。

对每一候选组合 $c=(\rho^Q,\theta,\omega)$，分别以SPT和ATC在全部30个实例上的平均目标值 $\overline F_b(c)$ 评价，并选取 $b^*(c)=\arg\min_{b\in\{\mathrm{SPT},\mathrm{ATC}\}}\overline F_b(c)$ 作为该组合的参考规则（相同时固定优先SPT），而非逐实例取两规则的最小值。记 $i$ 为校准实例，下列比例均基于同一参考规则的结果计算：
\begin{equation}\label{eq:service_calibration_statistics}
\begin{aligned}
  R_{\mathrm{WSF}}(c)&=\frac{\sum_i\mathrm{WSF}_{i,b^*(c)}(c)}{\sum_i F_{i,b^*(c)}(c)},\\
  R_{\mathrm{irr}}(c)&=\frac{\sum_i L_i(c)}{\sum_i\mathrm{WSF}_{i,b^*(c)}(c)},\\
  L_i(c)&=\sum_r w_{r2}\max\!\left\{0,Q_r-\sum_{j:e_j=r}q_j\right\}.
\end{aligned}
\end{equation}
其中 $L_i$ 是第3节所述\emph{最终检查点}的调度无关短缺下界，只在离线校准和结果分析时利用完整实例计算，决不进入在线策略的状态或动作。另令 $P_{\mathrm{early}}(c)$ 为30个实例中全部实体首次检查点发生正短缺的比例（分母为 $30R$），同样以 $b^*(c)$ 的排程计算。候选必须同时满足 $0.35\le R_{\mathrm{WSF}}(c)\le0.50$、$0.20\le P_{\mathrm{early}}(c)\le0.80$ 及 $R_{\mathrm{irr}}(c)\le0.40$；这些阈值用于事前排除阶段性要求几乎总能满足、几乎总是违约或最终不可避免短缺主导成本的退化场景，不是算法最优性的判据。分母为零时，相应候选视为不满足筛选要求。

在所有合格组合中，首先选取 $|R_{\mathrm{WSF}}(c)-0.425|$ 最小者；若并列，则依次取 $R_{\mathrm{irr}}(c)$ 更小者及 $\rho^Q,\theta,\omega$ 按升序字典序最小者。若没有合格组合，则报告候选网格和各筛选指标，不从本文方法或测试集表现反向选择参数，且须在正式规则生成和训练前另行声明并冻结修订协议。最终 $\rho^Q=\tbd[待校准]$、$\theta=\tbd[待校准]$、$\omega=\tbd[待校准]$；校准集的 $R_{\mathrm{WSF}}=\tbd[待校准]$，参考规则为 \tbd[待校准]，$P_{\mathrm{early}}=\tbd[待校准]$，$R_{\mathrm{irr}}=\tbd[待校准]$。对每个候选还报告 WSF 中超过已知下界的部分；该下界仅覆盖最终检查点，下界以上的短缺也可能包含无法由调度消除的早期损失，因而不能全部称为可改善成本。

预测快照由需求过程的已知统计参数构造，不读取已经抽样但尚未释放的作业。
预测采用齐次复合泊松模型，作为固定新到作业总数的实际到达机制的在线近似。它仅使用公开的到达率 $\lambda$、实体预期份额 $\pi_r$ 和服务数量分布的一二阶矩。各作业独立按 $\pi_r$ 归属实体，因此预测中实体 $r$ 的到达率设为 $\lambda_r=\lambda\pi_r$，未来长度为 $s\ge0$ 的区间内累计服务数量的均值、方差为
\begin{equation}\label{eq:entity_cp_forecast}
  \mu_r(t,t+s)=\lambda\pi_r s\,\mathbb{E}[q],\qquad
  v_r(t,t+s)=\lambda\pi_r s\,\mathbb{E}[q^2].
\end{equation}
预测覆盖至所有已公开检查点的最晚时刻，不读取未释放作业的实际实现。齐次模型可直接连续计算累计矩，无需人为设定预测桶宽；保守累计量仍按式~\eqref{eq:lambda} 在累计区间统一修正。对给定 $\kappa$，保守预测从零转为正的理论区间长度为 $s_{0r}=\kappa^2\mathbb{E}[q^2]/(\lambda\pi_r[\mathbb{E}[q]]^2)$，不同实体因 $\pi_r$ 不同而有不同的零值区间。风险系数的候选为 $\kappa\in\{0.25,0.50,0.75,1.00\}$。在固定的30个校准实例上，使用SPT与ATC产生的决策状态，统计每个尚未结算且存在正缺口的检查点在当前剩余区间内保守累计预测为零的比例 $P_0(\kappa)$，同时报告滚动累计预测对真实后续到达数量的误差和有正预测的时间窗口比例。为避免预测特征在绝大多数状态下恒为零或几乎始终为正，候选须满足 $0.20\le P_0(\kappa)\le0.80$；在合格候选中取 $\kappa$ 最大者以保留尽可能保守、但仍有状态区分力的预测；若无合格值则仅报告诊断并在规则生成与训练之前重新确认准则，不依靠PPO结果或测试数据挑选。最终 $\kappa=\tbd[待校准]$。上述比例仅为模型输入退化的事前诊断，不代表保守预测具有概率覆盖保证。

实际到达过程在第 $N_{\mathrm{new}}$ 个新作业到达后停止，复合泊松模型不考虑这一停止效应，也不使用 $N_{\mathrm{new}}$ 与已发生的到达次数，因此它是对未来到达的近似刻画，而非真实未来到达分布的精确描述：剩余计划到达数较多时，两者的均值接近；越接近到达结束，近似预测的均值越偏高，到达结束后仍给出正的预测量。本文采用这一近似，是为了使预测只依赖到达强度这类通用统计量，在计划到达总数不确定或无法事先给出的场景中同样适用。为说明近似带来的偏差大小，在校准实例上按到达进度分段报告预测累计量与实际到达量之间的绝对误差与平均偏差；二者之比只在实际到达量为正的时间窗上计算，实际无到达的时间窗单独统计其所占比例（\tbd[※]）。预测误差因此来自停止效应，以及到达时刻、依公开份额抽取的所属实体与服务数量的随机性。第4节的预测相关理论量均以这一近似预测为输入，其含义是``在该预测假设下''的风险信号。

中心场景取机器数 10、新到作业 50、目标利用率 0.80、交期紧度 1.0、服务权重倍率 $s_w=1$，共 70 个作业。策略只在中心场景训练，全部对比与消融实验在中心场景进行；在其余场景上的泛化检验见6.5节。

\subsubsection{数据划分与评价指标}

每个实例由随机种子及已冻结的场景参数唯一确定，同一个实例在各对比方法之间共用作业、实体份额、承诺与预测接口。实体份额、时限扰动、作业属性和到达时间分配独立随机数流；其中公开承诺不依赖未来作业样本。不同用途的种子区间互不相交，并由自动检验保证。各实例集的用途依次为：校准实例只用于确定服务参数及式~\eqref{eq:global-reward-scale}的全局奖励尺度（训练前固定，且不使用正式测试结果）；训练阶段每个回合都生成新实例，不重复使用；筛选实例（20 个）供 LLM 生成阶段与 GP 进化评价候选规则，并用于各规则库的选库；保留实例（20 个）用于复核所选规则在筛选实例之外的成绩，并检验不同生成批次之间的差异；验证实例（50 个）只用于学习率与折扣因子的选择，不参与训练；测试实例仅用于最终报告，中心场景 100 个，泛化场景各 50 个，压力测试另用 \tbd[30--50] 个独立实例。测试集事前计划取 100 个实例，并以配对差及其置信区间作为结论依据；是否需要更大样本须经独立功效分析或预实验确认，不以正式测试结果反向选择。

主指标为测试实例上的平均目标值 $F$，同时报告其分量 TT 与 WSF，数值越小越好。WSF 的调度无关下界 $\sum_r w_{rK_r}\underline{U}_r$（第3节）对所有方法相同，其在测试实例上占 WSF 的比例为 \tbd[※]，方法间的差异全部来自 WSF 中高于该下界的部分。涉及训练的方法各训练 3 个随机种子，评价时策略取概率最大的动作，使用训练结束时的最终模型，不按测试表现挑选检查点；方法间的统计比较在每个实例上取 3 个种子的平均值。主要推断依据是本文方法与各对比方法在同一实例上的配对差：给出自助法（bootstrap）95\% 置信区间与逐实例胜率，并做配对 Wilcoxon 符号秩检验，用 Holm 方法校正多重比较，显著性水平取 0.05。多种方法的整体排序以 Friedman 检验与 Nemenyi 事后比较~\cite{demsar2006}、以实例为区组的方差分析作为补充，不另作推断。

\subsubsection{对比方法}

对比方法分为七组（表~\ref{tab:methods}）。所有方法使用同一批实例；除滚动重排程遗传算法外，各方法共用同一非延迟派工机制和按机器编号确定触发机器的方式；规则类方法中规则只决定作业的优先次序。主对比中的学习型规则选择方法共用同一个规则库、同一检查点关系图编码器、同一真实奖励和相同的5000个训练回合，DDQN/A2C/Rainbow保留各自标准的标量价值更新，不使用特权未来信息和区间约束。Direct-Job PPO共用编码器、奖励与训练回合数，不使用规则库；用于隔离规则动作形式的主对照采用与SC--PPO相同的分解Critic、特权区间和 $\gamma=1$；另外保留普通标量Critic的Direct-Job PPO作为辅助对照，避免只比较两种均使用新价值机制的算法。

\begin{table}[t]
  \centering
  \caption{中心场景主对比的方法}\label{tab:methods}
  \small
  \begin{tabular}{@{}>{\raggedright}p{2.4cm}>{\raggedright}p{3.3cm}>{\raggedright\arraybackslash}p{9.0cm}@{}}
    \toprule
    组别 & 方法 & 说明 \\
    \midrule
    经典规则 & SPT、EDD、ATC、最小余量 & 最短加工时间、最早交货期、表观延迟成本、交期余量最小优先，单规则执行；最小余量即第5节的缺省排序 \\
    单规则基准 & Best-Single & 不学习，始终使用本文规则库中筛选实例上 $f_H$ 最小的一条规则（可能是 SPT 或 ATC）；用于区分规则质量与在线选择各自的贡献 \\
     & LLM-Best & 不学习，始终使用通过合法性检验与质量门槛的全部 LLM 候选规则中，筛选实例上 $f_H$ 最小的一条；用于与 GP-Best 比较单条规则的质量 \\
    自动规则生成 & GP-Best & 强类型遗传规划在相同三元组语法、特征与有效评价次数下演化出的最优单规则~\cite{xu2022gp,xu2024gp}，单规则执行 \\
    搜索类方法 & 滚动重排程遗传算法 & 每当有新作业到达，对所有已到达且尚未开始加工的工序重新搜索机器分配与加工顺序，其余时间按当前方案执行~\cite{luo2022}；两段式编码，种群 50，交叉概率 0.8，变异概率 0.2，每个重排程点限时 3 秒~\cite{lu2024,fu2025}；适应度为局部代理目标，见正文 \\
    学习型选择方法 & CRG--DDQN、CRG--A2C、CRG--Rainbow & 在同一规则库上以不同的强化学习算法训练同一检查点关系图策略：Double DQN~\cite{vanhasselt2016}；同步的 Advantage Actor--Critic，即 A3C 的同步版本~\cite{mnih2016}；Rainbow~\cite{hessel2018}，含双网络、对偶结构、优先回放、3 步回报、噪声网络与分布式价值 \\
    直接派工学习 & Direct-Job PPO & 不使用规则库，相同观测与编码器下直接在候选作业上以 masked-softmax 选作业~\cite{song2023} \\
     & Direct-Job PPO（标量Critic，辅助） & 仅把直接派工基线的区间分解Critic改为标准标量Critic；其余公开图、动作及训练预算与Direct-Job PPO一致 \\
    本文方法 & SC--PPO & 第5节方法 \\
    \bottomrule
  \end{tabular}
\end{table}

滚动重排程遗传算法代表在线场景下常见的搜索类做法。它与规则类方法的决策方式不同：它只在新作业到达时重新规划，但每次同时决定机器分配与加工顺序；规则类方法在每个可派工事件上决策，但只决定作业的优先次序，机器由编号顺序确定。两者的计算预算也不同，遗传算法每个重排程点搜索 3 秒，学习型方法每次决策只做一次前向计算（6.6 节）。遗传算法不使用预测，其适应度是只涉及已到达作业的局部代理目标：在当前方案下模拟已到达作业的完工时刻，累计其 TT；对每个尚未结算的检查点，只计已到达作业能够补足的那部分短缺，即以 $\min\{Q_{rk},\,Y_{rt}+\sum_{J_j\in\cJ_t,\,e_j=r}q_j\}$ 代替 $Q_{rk}$ 计算模拟短缺并加权求和。模拟与第3节的语义一致：已完工作业的数量已计入 $Y_{rt}$，已结算检查点的短缺已经确定，与方案无关，不进入适应度；正在加工的工序不可中断，按其实际完工时刻计入模拟；检查点恰在作业完工时刻到期时，该作业的数量计入。这样，尚未到达作业应承担的需求不计入适应度，避免把未来供给视为零而系统性高估服务短缺。该代理目标只用于每个重排程点的方案比较，方法的最终评价与其他方法相同，使用真实目标 $F$。因计算量大，该方法只在中心场景上评价，不参加泛化实验。其种群规模小于 Luo 等~\cite{luo2022} 的 200，并以限时代替固定的迭代次数，这一做法与 Lu 等~\cite{lu2024} 将元启发式限定在 3 秒内的设置一致。% 【待与代码核对：GA 适应度的现有实现】

GP-Best 是自动规则生成的标准基线。遗传规划超启发式是动态柔性作业车间中学习调度规则的常用方法~\cite{xu2022gp,xu2024gp,xu2026gp}。
% 【待核对：xu2026gp 指 ESWA 2026 的 GP with advanced diverse partner selection（DOI 10.1016/j.eswa.2026.132158），Claude 未能检索到原文，需本人确认后再保留】
为使 GP 与 LLM 的表达空间相同，GP 采用强类型遗传规划演化完整的三元组 $\langle\mathrm{Cond},\mathrm{Filter},\mathrm{Priority}\rangle$：Priority 为数值型表达式树，终端为表~\ref{tab:theory} 的 14 个特征与随机常数，函数集与 LLM 规则可用的运算一致（加、减、乘、保护除法、最大、最小）；Filter 与 Cond 为可选的布尔型分量，由特征（Cond 中为候选集合上该特征的最大值）与常数的比较经``与''``或''组合而成，深度上限 \tbd[※]。交叉与变异只在同类型分量之间进行：Priority 之间交换或替换子树，Filter 与 Cond 整体交换、改变比较符或阈值，并以一定概率增删 Filter 或 Cond。按此构造的规则在语法上总是合法的；运行时采用与 LLM 规则相同的除零、非有限值与条件失效约定（5.2.2 节），结果恒为常数、含非有限值或运行出错的规则同样剔除。GP 以筛选实例上的标准化目标 $f_H$ 为适应度。两者的主要公平性指标是有效候选评价次数，即在筛选实例上完整仿真的不同合法规则数：GP 的评价预算取 LLM 三个批次实际完成的有效候选评价次数，而非公式给出的上限 $3N_{\mathrm{gen}}N_c|\cI^{\mathrm{scr}}|$；重复个体通过缓存不重复计数。两者实际完成的有效评价次数、LLM 的调用次数、令牌与费用以及 GP 的进化耗时均如实报告（6.6 节）。GP 的种群规模、代数、交叉与变异概率为 \tbd[※]，进化结束后取筛选实例上 $f_H$ 最小的个体作为 GP-Best，单规则执行。
% 【待与代码核对：若 GP 实现只演化 Priority，需把上文改为“GP 只演化 Priority 分量，Cond 与 Filter 为空”，并删去“表达空间相同”的说法】

Direct-Job PPO 用于检验以规则为动作这一抽象本身的价值。它不使用启发式规则库，在相同的合法观测、触发机器、关系图编码器与奖励下，直接在触发机器的候选作业集合 $\mathcal J_n^{\mathrm{cand}}$ 上选择作业。记编码后候选作业 $j$ 的节点表示为 $h_j$、状态向量为 $z_n$（式~\eqref{eq:zn}），Actor 输出
\[
  u_j=\mathrm{MLP}_{\mathrm{a}}\bigl(h_j\,\Vert\,z_n\bigr),\qquad
  \pi(j\mid\cG_n)=\frac{\exp(u_j)}{\sum_{j'\in\mathcal J_n^{\mathrm{cand}}}\exp(u_{j'})},\quad j\in\mathcal J_n^{\mathrm{cand}},
\]
非候选作业不进入归一化，相当于 masked-softmax；主对照的训练Critic与SC--PPO相同，采用分解价值与区间认证；另提供普通标量Critic版本作为辅助对照，结果在6.3节单独列行，避免将其与区间匹配的主直接派工对照混为一谈。该对照受以图神经网络直接派工的深度强化学习方法启发~\cite{song2023}，但与其不同：Song 等的动作同时决定工序与机器，本对照固定触发机器，只选作业，以便与本文方法在相同的决策粒度下比较，因此它不是对该方法的完整复现。其训练回合数与PPO设置同本文方法；主对照采用 $\gamma=1$，学习率单独在6.2节选择。两种方法的比较同时报告目标值、环境交互步数与在线计算耗时。

\subsubsection{实现与训练设置}

规则生成使用千问（\tbd[模型版本号与调用时间]），通过 API 调用，不与其他 LLM 比较。独立生成 3 个批次，每个批次 \tbd[$N_{\mathrm{gen}}$] 轮、每轮 \tbd[$N_c$] 条候选，采样温度 \tbd[※]，候选在中心场景的筛选实例上评价；3 个批次至多需要 $3N_{\mathrm{gen}}N_c|\cI^{\mathrm{scr}}|$ 次仿真。3 个批次的候选合并为候选池后统一选库，规则库规模上限 $N_{\mathrm{lib}}=6$，混合库中含 SPT 与 ATC。3 个批次的候选总数为 \tbd[※]。作为参照，EoH~\cite{liu2024eoh} 在装箱问题上约用 2000 次查询，ReEvo~\cite{ye2024reevo} 限定 100 次评价。提示词、模型与采样参数见附录。

SC--PPO 的训练设置见表~\ref{tab:hyper}；其中学习率经6.2节的预实验确定。所有匹配学习对照使用式~\eqref{eq:global-reward-scale}在独立校准集上预先固定的同一 $F_{\mathrm{ref}}$，不再按每个训练实例重新计算。其余 PPO 超参数取该算法的常用值~\cite{schulman2017ppo,schulman2016gae}。计算环境为 \tbd[处理器、内存与加速器型号]。区间相关的训练特权未来信息只由Critic使用；在线推理阶段不使用。

\begin{table}[t]
  \centering
  \caption{SC--PPO 训练设置}\label{tab:hyper}
  \small
  \begin{tabular}{@{}>{\raggedright}p{2.8cm}>{\raggedright}p{3.4cm}>{\raggedright}p{2.6cm}>{\raggedright\arraybackslash}p{4.0cm}@{}}
    \toprule
    项目 & 取值 & 项目 & 取值 \\
    \midrule
    训练回合数 & 5000 & 每次更新回合数 $B$ & 32（共 157 次更新，最后一次为剩余的 8 个回合） \\
    每批数据更新轮次 $E$ & 4 & 学习率 & 见表~\ref{tab:lr} \\
    小批量大小 & \tbd[※] & 折扣因子 & $\gamma=1$（主SC--PPO） \\
    熵系数 & 0.02 线性降至 0.002 & GAE 参数 $\lambda_{\mathrm{GAE}}$ & 0.95 \\
    隐藏维度 & 64 & 奖励缩放 $c_s$ & 10 \\
    编码层数与头数 & 2 层，4 头（式~\eqref{eq:crg}） & 截断系数 & 0.2 \\
    价值损失系数 & 0.5 & 梯度裁剪 & 0.5 \\
    评价时动作 & 取概率最大的规则 & & \\
    \bottomrule
  \end{tabular}
\end{table}

三种学习型规则选择方法与Direct-Job PPO都与SC--PPO共用编码器结构、隐藏维度、训练回合数与评价时的贪心动作。SC--PPO和区间匹配的Direct-Job PPO固定 $\gamma=1$，各自预选学习率；CRG--DDQN、CRG--A2C、CRG--Rainbow各自验证学习率和折扣因子，保持各算法标准标量价值学习，不引入特权未来信息；三者主要参数见表~\ref{tab:rlbase}。

\begin{table}[t]
  \centering
  \caption{学习型对比方法的主要参数}\label{tab:rlbase}
  \small
  \begin{tabular}{@{}>{\raggedright}p{2.4cm}>{\raggedright\arraybackslash}p{11.6cm}@{}}
    \toprule
    方法 & 参数 \\
    \midrule
    CRG--DDQN & 经验回放容量 \tbd[※]；小批量 \tbd[※]；目标网络每 \tbd[※] 步同步；$\epsilon$-贪心探索，$\epsilon$ 由 \tbd[※] 线性降至 \tbd[※] \\
    CRG--A2C & 每次更新的回合数 \tbd[※]；优势估计 \tbd[$n$ 步回报或 GAE]；熵系数 \tbd[※]；价值损失系数 \tbd[※] \\
    CRG--Rainbow & 优先回放容量 \tbd[※]，优先级指数 $\alpha=$ \tbd[※]，重要性采样指数 $\beta$ 由 \tbd[※] 增至 1；3 步回报；分布式价值原子数 \tbd[※]，支撑区间 $[V_{\min},V_{\max}]=$ \tbd[※]；噪声网络初始标准差 \tbd[※]；目标网络每 \tbd[※] 步同步 \\
    \bottomrule
  \end{tabular}
\end{table}

%=====================================================================
\subsection{预实验：学习率、折扣因子与训练过程}

五类主学习型方法（SC--PPO、CRG--DDQN、CRG--A2C、CRG--Rainbow 与 Direct-Job PPO）共用同一编码器，但对学习率的敏感程度不同，为避免因参数未调而造成对比不公，每种方法分别在 $\{1\times10^{-4},3\times10^{-4},1\times10^{-3}\}$ 三档学习率中选择。每个配置训练 3 个随机种子、2000 个回合，训练实例与正式训练一样由训练种子区间逐回合新生成；训练过程中不使用验证实例更新参数，训练结束后在相同的 50 个验证实例上评估，以验证实例上平均目标值最低的一档作为正式实验的学习率（表~\ref{tab:lr}）。测试实例不参与任何参数选择。SC--PPO及区间匹配Direct-Job PPO在学习率预实验中固定 $\gamma=1$；其余三种学习算法的学习率预实验暂取 $\gamma=0.99$，随后单独选择折扣因子。

\begin{table}[t]
  \centering
  \caption{学习率预实验（验证实例，平均目标值$\pm$种子标准差）}\label{tab:lr}
  \small
  \begin{tabular}{@{}lcccc@{}}
    \toprule
    方法 & $1\times10^{-4}$ & $3\times10^{-4}$ & $1\times10^{-3}$ & 采用值 \\
    \midrule
    SC--PPO & \tbd[※] & \tbd[※] & \tbd[※] & \tbd[※] \\
    CRG--DDQN & \tbd[※] & \tbd[※] & \tbd[※] & \tbd[※] \\
    CRG--A2C & \tbd[※] & \tbd[※] & \tbd[※] & \tbd[※] \\
    CRG--Rainbow & \tbd[※] & \tbd[※] & \tbd[※] & \tbd[※] \\
    Direct-Job PPO & \tbd[※] & \tbd[※] & \tbd[※] & \tbd[※] \\
    \bottomrule
  \end{tabular}
\end{table}

\tbd[待汇总后填写：各方法对学习率的敏感程度，以及不同档位下种子间波动的差异。如保留熵系数、更新回合数、隐藏维度等 PPO 超参数的敏感性扫描，在此补充结果与结论。]

SC--PPO的区间Critic必须预测未经折扣的最终短缺，为保证Actor使用的TT与WSF保持与原总目标相同的时间尺度，SC--PPO及区间匹配Direct-Job PPO在所有训练与消融中固定 $\gamma=1$；不得对TT、WSF分支使用不同折扣因子。为保持其他学习算法的合理选择自由度，CRG--DDQN、CRG--A2C与CRG--Rainbow在选定学习率后各自比较 $\gamma\in\{0.99,0.999,1\}$，每档3种子、2000回合，并在相同验证实例上选择平均 $F$ 最低的一档（表~\ref{tab:gamma}）；这属于各基线的超参数验证，并不改变SC--PPO的固定目标口径。所有方法用未折扣真实 $F$ 评价。

\begin{table}[t]
  \centering
  \caption{折扣因子预实验（验证实例，平均目标值$\pm$种子标准差）}\label{tab:gamma}
  \small
  \begin{tabular}{@{}lcccc@{}}
    \toprule
    方法 & $\gamma=0.99$ & $\gamma=0.999$ & $\gamma=1$ & 采用值 \\
    \midrule
    SC--PPO & --- & --- & \textbf{固定为1} & 1 \\
    CRG--DDQN & \tbd[※] & \tbd[※] & \tbd[※] & \tbd[※] \\
    CRG--A2C & \tbd[※] & \tbd[※] & \tbd[※] & \tbd[※] \\
    CRG--Rainbow & \tbd[※] & \tbd[※] & \tbd[※] & \tbd[※] \\
    Direct-Job PPO & --- & --- & \textbf{固定为1} & 1 \\
    \bottomrule
  \end{tabular}
\end{table}

\begin{figure}[t]
  \centering
  \figph{fig6_1}{SC--PPO 在三档学习率下的训练曲线：验证集目标值随更新次数的变化，3 个种子的均值与范围}
  \caption{不同学习率下 SC--PPO 训练过程中验证集目标值的变化}\label{fig:lr}
\end{figure}

图~\ref{fig:lr} 将报告 SC--PPO 在三档学习率下的训练曲线。最终学习率下验证集目标值进入平稳区间的更新次数为 \tbd[※]，后期波动为 \tbd[※]；是否判定收敛由实测曲线决定。

%=====================================================================
\subsection{中心场景主对比}

表~\ref{tab:main} 给出中心场景（100 个测试实例）上各方法的平均目标值及分量，表~\ref{tab:tests} 给出本文方法相对各对比方法的配对检验，图~\ref{fig:forest} 为配对差森林图和平均序值比较。

\begin{table}[t]
  \centering
  \caption{中心场景主对比：测试实例平均目标值（学习型方法为种子均值$\pm$标准差）}\label{tab:main}
  \small
  \begin{tabular}{@{}lcccc@{}}
    \toprule
    方法 & $F$ & TT & WSF & 较最优经典规则的改进/\% \\
    \midrule
    SPT & \tbd[※] & \tbd[※] & \tbd[※] & \tbd[※] \\
    EDD & \tbd[※] & \tbd[※] & \tbd[※] & \tbd[※] \\
    ATC & \tbd[※] & \tbd[※] & \tbd[※] & \tbd[※] \\
    最小余量 & \tbd[※] & \tbd[※] & \tbd[※] & \tbd[※] \\
    库内最优单规则 & \tbd[※] & \tbd[※] & \tbd[※] & \tbd[※] \\
    LLM-Best & \tbd[※] & \tbd[※] & \tbd[※] & \tbd[※] \\
    GP-Best & \tbd[※] & \tbd[※] & \tbd[※] & \tbd[※] \\
    滚动重排程遗传算法 & \tbd[※] & \tbd[※] & \tbd[※] & \tbd[※] \\
    CRG--DDQN & \tbd[※] & \tbd[※] & \tbd[※] & \tbd[※] \\
    CRG--A2C & \tbd[※] & \tbd[※] & \tbd[※] & \tbd[※] \\
    CRG--Rainbow & \tbd[※] & \tbd[※] & \tbd[※] & \tbd[※] \\
    Direct-Job PPO & \tbd[※] & \tbd[※] & \tbd[※] & \tbd[※] \\
    Direct-Job PPO（标量Critic，辅助） & \tbd[※] & \tbd[※] & \tbd[※] & \tbd[※] \\
    SC--PPO（本文） & \tbd[※] & \tbd[※] & \tbd[※] & \tbd[※] \\
    \bottomrule
  \end{tabular}
\end{table}

\begin{table}[t]
  \centering
  \caption{本文方法相对各对比方法的配对检验}\label{tab:tests}
  \small
  \begin{tabular}{@{}lcccc@{}}
    \toprule
    对比方法 & 平均配对差 & 95\% 置信区间 & 胜率/\% & Wilcoxon $p$（Holm 校正） \\
    \midrule
    SPT & \tbd[※] & \tbd[※] & \tbd[※] & \tbd[※] \\
    EDD & \tbd[※] & \tbd[※] & \tbd[※] & \tbd[※] \\
    ATC & \tbd[※] & \tbd[※] & \tbd[※] & \tbd[※] \\
    最小余量 & \tbd[※] & \tbd[※] & \tbd[※] & \tbd[※] \\
    库内最优单规则 & \tbd[※] & \tbd[※] & \tbd[※] & \tbd[※] \\
    LLM-Best & \tbd[※] & \tbd[※] & \tbd[※] & \tbd[※] \\
    GP-Best & \tbd[※] & \tbd[※] & \tbd[※] & \tbd[※] \\
    滚动重排程遗传算法 & \tbd[※] & \tbd[※] & \tbd[※] & \tbd[※] \\
    CRG--DDQN & \tbd[※] & \tbd[※] & \tbd[※] & \tbd[※] \\
    CRG--A2C & \tbd[※] & \tbd[※] & \tbd[※] & \tbd[※] \\
    CRG--Rainbow & \tbd[※] & \tbd[※] & \tbd[※] & \tbd[※] \\
    Direct-Job PPO & \tbd[※] & \tbd[※] & \tbd[※] & \tbd[※] \\
    Direct-Job PPO（标量Critic，辅助） & \tbd[※] & \tbd[※] & \tbd[※] & \tbd[※] \\
    \bottomrule
  \end{tabular}
\end{table}

\begin{figure}[t]
  \centering
  \figph{fig6_2}{(a) 本文方法相对各对比方法的 F 配对差均值与 95\% 置信区间；(b) 各方法的平均序值与 Nemenyi 临界差}
  \caption{中心场景主对比：配对差与平均序值}\label{fig:forest}
\end{figure}

% 以下结论句待表~\ref{tab:main}、表~\ref{tab:tests} 数据确认后保留或修改
Friedman 检验的统计量为 \tbd[※]、$p=$ \tbd[※]，以实例为区组的方差分析的 $p=$ \tbd[※]；显著性及两类检验结论的一致性待实际结果确认。经典规则中表现最好的是 \tbd[规则名]，其平均目标值为 \tbd[※]。库内最优单规则来自 \tbd[LLM 规则/SPT/ATC]，相对最优经典规则的改进为 \tbd[※]\%。单条规则的质量以 LLM-Best 衡量：它相对最优经典规则的改进为 \tbd[※]\%，相对 GP-Best 的配对差为 \tbd[※]（95\% 置信区间 \tbd[※]），后者回答在相同语法、特征与有效评价次数下，LLM 与 GP 谁能生成更好的单条规则；库内最优单规则与 LLM-Best 可能相同，也可能因质量门槛、行为去重及边际选库而不同，须按最终规则编号核对。这一比较只涉及单条规则，不说明 LLM 规则库优于 GP 规则库。SC--PPO 相对库内最优单规则再改进 \tbd[※]\%（配对差 95\% 置信区间 \tbd[※]），在线选择带来的额外收益 \tbd[显著/不显著]。本文方法与 GP-Best 的配对差为 \tbd[※]（95\% 置信区间 \tbd[※]），这一差距同时包含规则来源与在线选择两部分。与滚动重排程遗传算法相比，本文方法的目标值 \tbd[更低/相当/更高]；两者的触发方式与决策权力不同，耗时对比见 6.6 节。与另外三种规则选择方法相比，SC--PPO 的平均目标值 \tbd[※]，种子间标准差 \tbd[※]。与 Direct-Job PPO 相比，配对差为 \tbd[※]，用以判断以规则为动作的抽象在相同观测与训练预算下是否有益。从分量看，本文方法的增益主要来自 \tbd[WSF/TT]。

%=====================================================================
\subsection{消融实验与规则选择行为}

消融实验在中心场景上进行，分别拆除完整方法中的一个部件，考察其作用。完整方法为含在线证书、训练区间约束与等价合并的SC--PPO，配合混合规则库（LLM 规则加 SPT、ATC）。各变体均训练 3 个种子，奖励参考尺度 $F_{\mathrm{ref}}$ 由独立校准实例上 SPT 与 ATC 的平均目标值确定（式~\eqref{eq:global-reward-scale}），对所有训练实例和规则库固定不变，确保训练奖励与原始期望总成本一致。

\subsubsection{规则库}

中心场景的混合规则库至多含 6 条规则，选库可因质量门槛、行为去重或边际增益不足而提前停止，实际为 \tbd[※] 条，其中 LLM 规则 \tbd[※] 条（表~\ref{tab:lib}）。表中同时给出各规则在保留实例上的成绩，用于检验选库结果是否依赖于筛选实例，以及每条 LLM 规则入库时的代理边际增益。混合库中 LLM 规则的数量在强化学习训练之前确认：若没有 LLM 规则进入混合库，则 PPO 的表现只能说明其在经典规则之间切换的作用，不能作为 LLM 规则贡献的证据；这一情形将作为实验发现报告，并据此检讨以逐实例最优成绩为代理目标的选库方式是否适合动态切换，而不是通过增加生成次数或放宽门槛凑足规则。

\begin{table}[t]
  \centering
  \caption{中心场景的规则库（筛选实例与保留实例上的成绩；行数按实际入选规则数调整）}\label{tab:lib}
  \small
  \setlength{\tabcolsep}{4pt}
  \begin{tabular}{@{}ccccccc>{\raggedright\arraybackslash}p{3.6cm}@{}}
    \toprule
    序号 & 来源 & 类别 & $f_H$（筛选） & $f_H$（保留） & 代理边际增益 & TT/WSF & 说明 \\
    \midrule
    1 & \tbd[※] & \tbd[※] & \tbd[※] & \tbd[※] & \tbd[※] & \tbd[※] & \tbd[※] \\
    2 & \tbd[※] & \tbd[※] & \tbd[※] & \tbd[※] & \tbd[※] & \tbd[※] & \tbd[※] \\
    3 & \tbd[※] & \tbd[※] & \tbd[※] & \tbd[※] & \tbd[※] & \tbd[※] & \tbd[※] \\
    4 & \tbd[※] & \tbd[※] & \tbd[※] & \tbd[※] & \tbd[※] & \tbd[※] & \tbd[※] \\
    5 & \tbd[※] & \tbd[※] & \tbd[※] & \tbd[※] & \tbd[※] & \tbd[※] & \tbd[※] \\
    6 & \tbd[※] & \tbd[※] & \tbd[※] & \tbd[※] & \tbd[※] & \tbd[※] & \tbd[※] \\
    \bottomrule
  \end{tabular}
\end{table}

表~\ref{tab:abl_lib} 在相同的网络和训练预算下比较不同的规则库。比较关系如下：与纯 LLM 库相比，检验把 SPT、ATC 并入库中是否必要；与原始可见量特征库相比，检验第4节理论特征对规则生成的作用；与手工规则库和随机规则库相比，检验 LLM 规则相对人工编写和随机采样的优势；与经典规则库相比，检验加入 LLM 规则本身的收益。

原始可见量特征库用于检验理论特征对规则生成的作用：LLM 可用的规则特征由表~\ref{tab:theory} 的 14 个特征换为 7 个未经第4节推导的原始可见量，即作业服务数量（按其期望缩放）、所属实体最早的缺口为正的未结算检查点的缺口比例与剩余时间，以及 4 个基本调度特征；生成批次、轮数、每轮候选数与选库流程均与本文相同。该变体只改变规则可用的特征，PPO 的状态表示仍为表~\ref{tab:feat} 的完整特征，因此它与本文的差异来自规则库本身，而非状态编码。两者的规则特征维数与信息量均不同，这一比较检验的是理论特征集作为整体对规则生成的作用，不区分单个特征、特征维数或信息量的影响，也不能单独归因于理论推导本身。

手工规则库为依据第4节理论量人工编写的 \tbd[※] 条规则。所有规则来源（LLM、原始可见量 LLM、手工、随机）采用同一套筛选协议：在相同的筛选实例上做合法性检验、质量门槛（$f_H<1$）、行为去重与边际增益选库，再并入 SPT 与 ATC。随机规则库按同样的规则结构与特征集随机采样，候选总数与 LLM 三个批次的候选总数相同。各来源的质量门槛通过率单独报告（表~\ref{tab:abl_lib}）；若某一来源通过门槛的规则不足以填满规则库，则如实使用较小的规则库，不为凑足规模而放宽门槛，规则库的实际规模同时列于表中。经典规则库只含 4 条规则，其比较同时包含规则来源与规则数量两方面的差异。

\begin{table}[t]
  \centering
  \caption{规则库消融（中心场景，SC--PPO 选择）}\label{tab:abl_lib}
  \small
  \setlength{\tabcolsep}{4pt}
  \begin{tabular}{@{}>{\raggedright}p{4.4cm}ccccccc@{}}
    \toprule
    规则库 & 通过率/\% & 库规模 & $F$ & TT & WSF & 配对差（相对本文） & 95\% 置信区间 \\
    \midrule
    本文：LLM 规则 + SPT、ATC & \tbd[※] & \tbd[※] & \tbd[※] & \tbd[※] & \tbd[※] & --- & --- \\
    纯 LLM 规则 & \tbd[※] & \tbd[※] & \tbd[※] & \tbd[※] & \tbd[※] & \tbd[※] & \tbd[※] \\
    LLM 规则（原始可见量特征）+ SPT、ATC & \tbd[※] & \tbd[※] & \tbd[※] & \tbd[※] & \tbd[※] & \tbd[※] & \tbd[※] \\
    手工规则 + SPT、ATC & \tbd[※] & \tbd[※] & \tbd[※] & \tbd[※] & \tbd[※] & \tbd[※] & \tbd[※] \\
    随机规则 + SPT、ATC & \tbd[※] & \tbd[※] & \tbd[※] & \tbd[※] & \tbd[※] & \tbd[※] & \tbd[※] \\
    经典规则（SPT、EDD、ATC、最小余量） & --- & 4 & \tbd[※] & \tbd[※] & \tbd[※] & \tbd[※] & \tbd[※] \\
    \bottomrule
  \end{tabular}
\end{table}

\tbd[结论句：理论特征相对原始可见量的增益；LLM 规则相对随机、手工与经典规则库的差距及显著性；纯 LLM 库与混合库的差异。待数据确认后填写。]

规则库、状态表示与学习机制消融的结论按四条证据链组织：其一，理论特征对规则生成的增量（本文规则库对原始可见量特征库）；其二，LLM 规则生成相对其他生成方式的增量（单规则层面为 LLM-Best 对 GP-Best，规则库层面为本文对手工、随机与经典规则库），以及相应的生成成本；其三，按状态动态选择规则的增量（本文对库内最优单规则，以及 6.4.3 节的频率匹配选择与随机选择）。其四，区间认证的独立增量（表~\ref{tab:critic}的C4对C3，结合图~\ref{fig:cert}的证书紧度），并将Critic估计误差与最终调度目标分开报告。任何一条证据链若配对差的置信区间包含 0，即如实报告该部分收益未得到统计支持，不作超出数据的解读。

\subsubsection{状态表示与学习机制}

表~\ref{tab:abl_net} 在相同规则库和相同训练预算下检验状态表示各部分的作用。去除图结构的变体以作业特征的均值与最大值、检查点特征的均值和全局特征拼接后经多层感知机输入 PPO。去除检查点节点的变体在图中删去检查点节点及其关联边，作业特征中的服务信息保留，因此它检验的是显式检查点节点的作用，而非服务信息整体的作用。该变体仍使用C4分解Critic：对各检查点头，仅以其合法可见12维原始检查点特征经独立线性投影形成不含检查点关系消息的辅助表示，连同与完整方法相同的在线/训练特权区间及私有摘要输入Critic；Actor不读取这些辅助检查点表示，以确保\emph{去节点}对照不会经Critic私有通道还原Actor的完整检查点图。去除预测特征的变体把状态与规则中全部依赖预测快照的量置零，包括作业的预测不足特征，检查点节点的预测不足比例、预测到达余量与预测余量标志，以及全局的预测不足标志（检查点节点的未来需求比例只依赖可见作业，予以保留）；规则库不变，规则执行时读取到置零后的值。去除容量特征的变体以同样方式置零容量压力、容量标志及对应的检查点与全局特征。不依赖网络的选择机制对照（随机选择与频率匹配选择）见 6.4.3 节。

\begin{table}[t]
  \centering
  \caption{状态表示消融（中心场景，混合规则库）}\label{tab:abl_net}
  \small
  \begin{tabular}{@{}lccccc@{}}
    \toprule
    变体 & $F$ & TT & WSF & 配对差（相对本文） & 95\% 置信区间 \\
    \midrule
    完整 SC--PPO（本文） & \tbd[※] & \tbd[※] & \tbd[※] & --- & --- \\
    去除图结构（MLP--PPO） & \tbd[※] & \tbd[※] & \tbd[※] & \tbd[※] & \tbd[※] \\
    去除检查点节点 & \tbd[※] & \tbd[※] & \tbd[※] & \tbd[※] & \tbd[※] \\
    不区分可达与不可达关系 & \tbd[※] & \tbd[※] & \tbd[※] & \tbd[※] & \tbd[※] \\
    去除在线证书区间两维 & \tbd[※] & \tbd[※] & \tbd[※] & \tbd[※] & \tbd[※] \\
    去除预测特征 & \tbd[※] & \tbd[※] & \tbd[※] & \tbd[※] & \tbd[※] \\
    去除容量特征 & \tbd[※] & \tbd[※] & \tbd[※] & \tbd[※] & \tbd[※] \\
    \bottomrule
  \end{tabular}
\end{table}

\paragraph{区间认证Critic的独立贡献} 为分离价值分解、训练特权信息和数学区间约束的作用，在\emph{完全相同的公开图编码器、动作空间、LLM规则库、等价合并开关、训练种子与计算预算}下构造四组Critic。C1为标准标量Critic；C2为仅使用合法在线信息的分解Critic；C3为使用与C4完全相同的检查点节点编码、在线区间端点、训练特权外生特征以及特权区间端点作为输入，但输出头采用非负无区间参数化的分解Critic；C4为式~\eqref{eq:cert-critic}的原冻结标量区间认证分解Critic（原始完整方法基线）。全部采用统一 $\gamma=1$、相同奖励及相同GAE参数；C3与C4控制相同的特权信息访问权限，以防把特权信息本身的收益误归因于区间约束。四组使用相同的未折扣真实成本标签：C1以完整蒙特卡洛剩余总成本监督标量头，C2--C4按式~\eqref{eq:cert-critic-loss-c4}的逐检查点直接求和平方损失分项监督，三组采用相同的价值损失系数 $c_v$、梯度裁剪与蒙特卡洛标签；C3与C4仅在最后的区间限制方式不同，不额外加入TD目标裁剪。\emph{主要识别对照为C4--C3}，而非仅比较C4与C1。

\begin{table}[t]
  \centering
  \caption{区间认证Critic的递进消融（中心场景，相同公开编码器与规则库）}\label{tab:critic}
  \small
  \begin{tabular}{@{}p{3.3cm}ccccc@{}}
    \toprule
    方法 & 特权信息 & 硬区间约束 & $F$ & Critic误差 & 收敛回合/步 \\
    \midrule
    C1：标量Critic & 否 & 否 & \tbd[※] & \tbd[※] & \tbd[※] \\
    C2：分解Critic & 否 & 否 & \tbd[※] & \tbd[※] & \tbd[※] \\
    C3：特权分解Critic & 是 & 否 & \tbd[※] & \tbd[※] & \tbd[※] \\
    C4：标量区间认证Critic（基线） & 是 & 是 & \tbd[※] & \tbd[※] & \tbd[※] \\
    \bottomrule
  \end{tabular}
\end{table}


\paragraph{CC--DVL的分布价值学习消融}
在现有C1--C4标量Critic消融基础上，增加D1和D2两组分位数Critic，用于分离分布学习与确定性证书限制的贡献。D1为\emph{不施加硬证书输出约束}的有序分位数Critic，采用非负、无硬上界的有序参数化；D2为式~\eqref{eq:cc-quantile}的\emph{证书约束有序分位数Critic}，即本节待验证的CC--DVL候选主方法。D1和D2均使用式~\eqref{eq:cc-quantile-loss}所定义的原尺度pinball损失形式及相同中点分位数网格；D1仅将其中的预测量替换成其无硬区间输出。式~\eqref{eq:cc-pinball-equivalence}已证明乘回证书宽度的归一化pinball形式与原尺度损失恒等，因此不把两种写法视为不同的训练机制，也不设置数学上重复的对照组。

为识别机制增量，D1、D2与C4应固定相同的公开状态图、Actor、动作集合、LLM规则库、等价合并方式、训练期特权字段（包括两种证书端点）、训练/测试种子、奖励与 $\gamma=1$ 的GAE协议，且在同一全局 $F_{\mathrm{ref}}$ 下使用相同的环境交互次数与优化器更新预算。D1虽可\emph{读取}特权区间端点，却不得用它们对输出进行截断；D1与D2的唯一结构性差异为是否把硬证书嵌入分位数输出。D1、D2使用相同真实结算短缺标签，不得对D2额外提供数据。主要对照为\emph{D2--D1}（硬区间约束的分布学习增量）及\emph{D2--C4}（分布式与标量证书Critic的增量）；C4与D2使用不同种类的监督损失，应报告损失系数、参数量、墙钟训练时长和实际梯度计算量，不将单纯的模型容量差异归结为证书收益。分位数数目 $K$ 与损失系数 $\alpha_Q$ 仅使用事先划分的验证集选择，保持预设且相当的超参数搜索预算，不依据测试轨迹调整。
\begin{table}[t]
\centering
\caption{CC--DVL候选价值学习机制的匹配对照（数据待实验填入）}\label{tab:ccdvl_ablation}
\begin{tabular}{lccc}
\hline
模型 & 短缺分位数 & 硬证书 & $F$（测试）\\
\hline
C4：冻结标量证书Critic & 否 & 是 & \tbd[待测]\\
D1：无硬证书分位数Critic & 是 & 否 & \tbd[待测]\\
D2：CC--DVL证书分位数Critic & 是 & 是 & \tbd[待测]\\
\hline
\end{tabular}
\end{table}
除真实$F$、TT、WSF外，报告检查点pinball误差、经验分位数校准偏差、均值短缺MAE、Critic总价值误差、优势估计方差、训练稳定性及在线推理耗时；证书真实标签包含率单独验证且应为100\%。由于训练期特权信息可能使条件不确定性收缩，按有效证书宽度与条件短缺波动程度分层分析分位数是否退化。若D2相对C4没有可信的独立增量，应如实报告，并保留原标量C4作为论文主方法候选，不预先宣称分布学习带来性能改善。

\paragraph{证书有效性与紧度} 对中心场景保留实例及测试实例记录每个活动检查点的在线/特权证书包含率、归一化宽度（两个版本统一使用式~\eqref{eq:cert-vis}的容量收紧可见供给上界） $\bigl(H^x_{rk}-L^x_{rk}\bigr)/Q_{rk}$、到截止时间的相对距离及区间历史收缩过程。严格界的包含率应为100\%；违反时必须修复生成器、证明或代码而不得当作普通预测误差。图~\ref{fig:cert} 给出在不同时间阶段的区间宽度分布，区分在线版与训练特权版，不使用测试结果调整边界。另记录C3/C4的TT预测误差、检查点短缺MAE、总价值估计误差和PPO训练曲线；为避免特权头隐性泄漏，验证/测试时仅以公开Actor派工并事后对固定轨迹评估Critic预测误差，不用特权预测参与选动作，以及最终 $F$、TT、WSF，避免仅以Critic误差代替调度收益。

\begin{figure}[t]
 \centering
 \figph{fig6_cert_v2}{活动检查点短缺证书区间宽度：在线/训练特权两组曲线，按相对截止时间与训练决策阶段分组}
 \caption{服务证书的有效性及区间紧度诊断}\label{fig:cert}
\end{figure}

\paragraph{等价合并消融} 在固定图、规则库与C4的情况下，分别开启/关闭式~\eqref{eq:equiv}的规则动作等价合并。报告 $F$、TT、WSF、有效分歧率、Actor有效学习步数及额外枚举规则的在线耗时。关闭合并时按原始规则索引计算PPO概率、比率和熵；开启时这些量统一在动作等价类上计算。两者均保留所有环境转移。

\begin{figure}[t]
  \centering
  \figph{fig6_3}{规则库、状态表示与区间Critic消融的配对差，按组显示95\%置信区间}
  \caption{主要消融实验的配对差}\label{fig:abl}
\end{figure}

\tbd[结论句：去除图结构、检查点节点、预测特征、容量特征后 F 的变化及显著性。待数据确认后填写。]

\subsubsection{规则选择机制与行为分析}

本小节考察 PPO 学到的选择是否依赖状态，以及规则切换是否真正改变派工结果。三者均使用完整方法训练所得的模型与混合规则库，不另行训练。

\paragraph{选择机制对照} 表~\ref{tab:select} 在同一规则库上比较四种选择方式。随机选择在每个决策点从规则库中均匀随机选一条规则，作为不依赖状态与网络的下限。频率匹配选择（Freq-Matched）先在完整 SC--PPO 的验证轨迹上，计算每条原始规则的平均策略概率 $\widehat p_h=|\mathcal S_{\mathrm{val}}|^{-1}\sum_{s\in\mathcal S_{\mathrm{val}}}\pi_\theta(H_h\mid s)$，再在每个决策点独立按 $\widehat p_h$ 抽样规则，不读取当前状态。由于主方法启用同派工动作等价合并，此处匹配的是原始规则的\emph{边际概率质量}，而不是由执行作业唯一恢复出的规则次数；更不保证新策略改变轨迹后实际使用频率不变。若本文方法同时优于库内最优单规则与频率匹配选择，可作为状态条件选择有益的经验依据，但并非严格的因果证明。

\begin{table}[t]
  \centering
  \caption{规则选择机制对照（中心场景，混合规则库）}\label{tab:select}
  \small
  \begin{tabular}{@{}lccccc@{}}
    \toprule
    选择方式 & $F$ & TT & WSF & 配对差（相对本文） & 95\% 置信区间 \\
    \midrule
    SC--PPO（本文，按状态选择） & \tbd[※] & \tbd[※] & \tbd[※] & --- & --- \\
    库内最优单规则（固定一条） & \tbd[※] & \tbd[※] & \tbd[※] & \tbd[※] & \tbd[※] \\
    频率匹配选择（Freq-Matched） & \tbd[※] & \tbd[※] & \tbd[※] & \tbd[※] & \tbd[※] \\
    随机选择 & \tbd[※] & \tbd[※] & \tbd[※] & \tbd[※] & \tbd[※] \\
    \bottomrule
  \end{tabular}
\end{table}

\paragraph{规则的有效分歧} 规则使用比例的变化并不等于派工结果的变化：若两条规则在某状态下选出同一作业，切换规则在执行上没有区别。因此在完整方法的测试轨迹上，对候选作业不少于 2 个的决策状态集合 $\mathcal{S}^{\geq2}$ 计算规则库的有效分歧率
\[
  D_{\mathrm{eff}}=\frac{\bigl|\{s\in\mathcal{S}^{\geq2}:\ \text{库内至少两条规则在 } s \text{ 选出不同作业}\}\bigr|}{|\mathcal{S}^{\geq2}|},
\]
同时报告 PPO 实际派工的作业与 SPT、ATC 及库内最优单规则在同一状态下所选作业不同的比例，以及库内每对规则在 $\mathcal{S}^{\geq2}$ 上首选作业不同的比例（$1-\mathrm{Sim}$，式~\eqref{eq:sim}）。这些量只依赖日志，不需要额外训练，用于判断规则库是否提供了值得学习的不同调度行为（\tbd[※]）。

\paragraph{状态与选择的关联} 为考察 PPO 学到的选择是否与状态相关，在中心场景测试实例上记录原始规则策略概率分布，包括按各决策状态的平均概率质量计算的规则占比和 LLM 规则概率占比。为绘制相邻决策的规则切换轨迹，仅在诊断日志中将每个已执行等价类内原始概率最大的规则作为代表（同分取编号最小者）；该代表不参与动作采样，也不能解释成合并后执行动作唯一对应的规则。按检查点风险、交期风险与容量压力对决策状态分组：三类风险分别取决策点上目标检查点的预测不足特征、候选作业的最大拖期暴露与容量压力特征，各按验证实例上的上下四分位数分为低、中、高三档，分档阈值在测试前确定，不依据测试结果回调。图~\ref{fig:behavior}(a) 以风险分组为行、入库规则为列，单元格为该组状态下各原始规则的平均策略概率质量；图~\ref{fig:behavior}(b) 给出一个典型测试实例的检查点累计完成曲线与规则切换轨迹。行为分析只描述选择与状态之间的关联，不作为因果证据，其解释需结合表~\ref{tab:select} 的选择机制对照与上述有效分歧率。

\begin{figure}[t]
  \centering
  \figph{fig6_6}{(a) 规则选择条件概率热图：行为按验证实例四分位数划分的检查点风险、交期风险与容量压力分组，列为入库规则，单元格为平均选择比例；(b) 典型测试实例的检查点累计完成曲线与规则切换轨迹}
  \caption{规则选择行为分析}\label{fig:behavior}
\end{figure}

\tbd[结论句：本文相对频率匹配与随机选择的配对差；有效分歧率与 PPO 派工偏离经典规则的比例；LLM 规则的平均概率质量占比与诊断代表规则的切换率；高风险状态下选择比例的变化。待数据确认后填写。]

%=====================================================================
\subsection{未见场景上的泛化与压力测试}

策略只在中心场景训练。为检验其在其他车间条件下的表现，将训练好的模型不再重新训练，直接用于另外 9 个场景，每个场景 50 个测试实例；规则库也沿用中心场景的规则库。完整的 $3^4=81$ 个因子组合数量过大；文献中的全因子设计常见 27 至 243 个配置~\cite{zheng2025,lei2024}，结果表难以阅读。本文采用 $L_9(3^4)$ 正交表选取 9 个场景（表~\ref{tab:l9def}），四个因子的每一档恰好各出现 3 次，可以比较各因子的主效应，但无法分离因子间的交互作用。每个实体的检查点数量（$K_r=2$）、需求份额生成分布 $a_r\sim U[0.8,1.2]$、预测协议与中心场景校准的基础服务参数在全部场景中保持不变；总作业数固定为70，实体数均为9，检查点总数均为18，但每个实例的 $\pi_r$ 独立抽样。服务权重倍率只对基础权重 $w_{r1}=\omega$、$w_{r2}=3\omega$ 施加共同缩放，即 $w_{r1}=s_w\omega$、$w_{r2}=3s_w\omega$；第3节目标仍为 $F=TT+WSF$。因此该L9设计混合了车间条件泛化和目标权衡敏感性，不能将 $s_w$ 的主效应解释为加工规模变化；$L_9$ 只允许讨论主效应，不能识别交互作用。

\begin{table}[t]
  \centering
  \caption{泛化场景（$L_9$ 正交表；实测利用率及SPT的TT占比待新实例协议校准完成后重算）}\label{tab:l9def}
  % 实测利用率和 TT/F 均必须依据新生成器、与正式测试种子重新统计；旧数据不可沿用。
  % 论文跨章同步：sec3的公开先验应包含pi_r；sec4的lambda_r=lambda*pi_r；sec5维度保持不变。
  \small
  \begin{tabular}{@{}lccccc>{\centering\arraybackslash}p{2.4cm}@{}}
    \toprule
    场景 & 机器数 & 服务倍率 $s_w$ & 目标利用率 & 交期紧度 & 实测利用率 & SPT 的 TT 占 $F$/\% \\
    \midrule
    中心 & 10 & 1 & 0.80 & 1.0 & \tbd[重算] & \tbd[重算] \\
    1 & 5 & 0.2 & 0.65 & 0.5 & \tbd[重算] & \tbd[重算] \\
    2 & 5 & 1 & 0.80 & 1.0 & \tbd[重算] & \tbd[重算] \\
    3 & 5 & 5 & 0.95 & 1.5 & \tbd[重算] & \tbd[重算] \\
    4 & 10 & 0.2 & 0.80 & 1.5 & \tbd[重算] & \tbd[重算] \\
    5 & 10 & 1 & 0.95 & 0.5 & \tbd[重算] & \tbd[重算] \\
    6 & 10 & 5 & 0.65 & 1.0 & \tbd[重算] & \tbd[重算] \\
    7 & 20 & 0.2 & 0.95 & 1.0 & \tbd[重算] & \tbd[重算] \\
    8 & 20 & 1 & 0.65 & 1.5 & \tbd[重算] & \tbd[重算] \\
    9 & 20 & 5 & 0.80 & 0.5 & \tbd[重算] & \tbd[重算] \\
    \bottomrule
  \end{tabular}
\end{table}

表~\ref{tab:l9res} 给出各场景的结果。参与泛化评价的方法包括 SPT、ATC、库内最优单规则、GP-Best、三种学习型规则选择方法、Direct-Job PPO 与本文方法，其中库内最优单规则与 GP-Best 在中心场景的筛选实例上确定后固定使用，学习型方法均使用中心场景训练的模型。如果本文方法在其他场景上稳定优于这条固定规则，则说明在线选择能够随状态变化调整所用规则，这是固定执行一条规则无法做到的。

\begin{table}[t]
  \centering
  \caption{泛化场景上的平均目标值 $F$}\label{tab:l9res}
  \footnotesize
  \setlength{\tabcolsep}{3pt}
  \begin{tabular}{@{}lcccccccccc@{}}
    \toprule
    场景 & SPT & ATC & 最优单规则 & GP-Best & DDQN & A2C & Rainbow & 直接派工 & 本文 & 改进/\% \\
    \midrule
    1 & \tbd[※] & \tbd[※] & \tbd[※] & \tbd[※] & \tbd[※] & \tbd[※] & \tbd[※] & \tbd[※] & \tbd[※] & \tbd[※] \\
    2 & \tbd[※] & \tbd[※] & \tbd[※] & \tbd[※] & \tbd[※] & \tbd[※] & \tbd[※] & \tbd[※] & \tbd[※] & \tbd[※] \\
    3 & \tbd[※] & \tbd[※] & \tbd[※] & \tbd[※] & \tbd[※] & \tbd[※] & \tbd[※] & \tbd[※] & \tbd[※] & \tbd[※] \\
    4 & \tbd[※] & \tbd[※] & \tbd[※] & \tbd[※] & \tbd[※] & \tbd[※] & \tbd[※] & \tbd[※] & \tbd[※] & \tbd[※] \\
    5 & \tbd[※] & \tbd[※] & \tbd[※] & \tbd[※] & \tbd[※] & \tbd[※] & \tbd[※] & \tbd[※] & \tbd[※] & \tbd[※] \\
    6 & \tbd[※] & \tbd[※] & \tbd[※] & \tbd[※] & \tbd[※] & \tbd[※] & \tbd[※] & \tbd[※] & \tbd[※] & \tbd[※] \\
    7 & \tbd[※] & \tbd[※] & \tbd[※] & \tbd[※] & \tbd[※] & \tbd[※] & \tbd[※] & \tbd[※] & \tbd[※] & \tbd[※] \\
    8 & \tbd[※] & \tbd[※] & \tbd[※] & \tbd[※] & \tbd[※] & \tbd[※] & \tbd[※] & \tbd[※] & \tbd[※] & \tbd[※] \\
    9 & \tbd[※] & \tbd[※] & \tbd[※] & \tbd[※] & \tbd[※] & \tbd[※] & \tbd[※] & \tbd[※] & \tbd[※] & \tbd[※] \\
    \midrule
    平均 & \tbd[※] & \tbd[※] & \tbd[※] & \tbd[※] & \tbd[※] & \tbd[※] & \tbd[※] & \tbd[※] & \tbd[※] & \tbd[※] \\
    \bottomrule
  \end{tabular}
\end{table}

\begin{figure}[t]
  \centering
  \figph{fig6_4}{四个因子每一档上本文方法相对 SPT 与 ATC 中较优者的改进百分比（主效应图），每一档取对应 3 个场景的平均}
  \caption{各因子水平上的主效应}\label{fig:effects}
\end{figure}

表中改进百分比以 SPT 与 ATC 中较优者为基准。本文方法在 9 个场景中的 \tbd[※] 个场景上优于基准，改进幅度从 \tbd[※]\% 到 \tbd[※]\%。图~\ref{fig:effects} 显示，\tbd[结论句：优势随机器数、利用率或交期紧度变化的趋势，以及优势消失的场景及原因。待数据确认后填写。]

上述泛化不改变服务结构与预测协议。为在中心场景上单独考察这两方面，另做两类小规模压力测试，均使用中心场景训练的模型，不重新训练，测试实例为 \tbd[30--50] 个配对实例，种子区间与其他用途互不相交。测试只用于报告方法对训练外条件的敏感程度，不用于回调任何参数或挑选测试条件。

其一，每个实体的检查点数量 $K_r\in\{1,2,3\}$。三档保持 $Q_r$、最终截止时间与计划服务时间窗不变，$K_r=1$ 时只保留最终检查点，$K_r=3$ 时在两个现有检查点之间增加一个，其时刻与累计需求比例为 \tbd[※]；三档的检查点权重总和保持相同，避免成本尺度随检查点数量自然升高而被误读为方法效果。$K_r=1$ 与 $K_r=3$ 均属训练外分布。

其二，预测失配。仅将策略与规则所用的复合泊松总预测到达率设为 $0.7\lambda$、$1.0\lambda$、$1.3\lambda$，即实体预测率统一替换为 $0.7\lambda\pi_r$、$1.0\lambda\pi_r$、$1.3\lambda\pi_r$，公开份额、服务承诺及实际抽样的到达轨迹均保持不变，以检验预测偏低或偏高时方法的稳健性。

参与压力测试的方法为 SPT、ATC、库内最优单规则、Direct-Job PPO 与本文方法，结果见表~\ref{tab:stress}。不同 $K_r$ 下检查点的时刻与加权需求结构不同，问题本身的难度随之改变，因此不同档位之间的绝对目标值不能直接比较：$K_r=3$ 时 $F$ 较高，可能只说明问题变难，而不说明方法退化。结论只依据同一档位、同一批实例上本文方法相对竞争方法的配对比较，表中给出本文方法相对库内最优单规则与 Direct-Job PPO 的改进百分比及配对差的 95\% 置信区间；预测失配测试同样按档位内的配对比较解读。

\begin{table}[t]
  \centering
  \caption{中心场景压力测试（配对实例）：各档位内的平均目标值与本文方法的相对改进}\label{tab:stress}
  \footnotesize
  \setlength{\tabcolsep}{3pt}
  \begin{tabular}{@{}llccccccc@{}}
    \toprule
     & & \multicolumn{5}{c}{平均目标值 $F$} & \multicolumn{2}{c}{本文的改进/\%（95\% 置信区间）} \\
    \cmidrule(lr){3-7}\cmidrule(l){8-9}
    测试 & 设置 & SPT & ATC & 最优单规则 & Direct-Job & 本文 & 相对最优单规则 & 相对 Direct-Job \\
    \midrule
    检查点数量 & $K_r=1$ & \tbd[※] & \tbd[※] & \tbd[※] & \tbd[※] & \tbd[※] & \tbd[※] & \tbd[※] \\
     & $K_r=2$（训练设定） & \tbd[※] & \tbd[※] & \tbd[※] & \tbd[※] & \tbd[※] & \tbd[※] & \tbd[※] \\
     & $K_r=3$ & \tbd[※] & \tbd[※] & \tbd[※] & \tbd[※] & \tbd[※] & \tbd[※] & \tbd[※] \\
    \midrule
    预测到达率 & $0.7\lambda$ & --- & --- & \tbd[※] & \tbd[※] & \tbd[※] & \tbd[※] & \tbd[※] \\
     & $1.0\lambda$（训练设定） & --- & --- & \tbd[※] & \tbd[※] & \tbd[※] & \tbd[※] & \tbd[※] \\
     & $1.3\lambda$ & --- & --- & \tbd[※] & \tbd[※] & \tbd[※] & \tbd[※] & \tbd[※] \\
    \bottomrule
  \end{tabular}
\end{table}

预测失配测试中 SPT 与 ATC 不使用预测，结果与 $1.0\lambda$ 相同，不重复列出；库内最优单规则若不含预测特征，同样不受影响。\tbd[结论句：各档位内本文方法相对库内最优单规则与 Direct-Job PPO 的改进及其置信区间如何随检查点数量与预测失配变化。待数据确认后填写。]

%=====================================================================
\subsection{收敛过程与计算开销}

图~\ref{fig:converge} 比较五种学习型方法的训练过程，浅色线为各种子的原始曲线，深色线为 3 个种子均值的滑动平均，阴影为 3 个种子的取值范围。\tbd[结论句：各方法达到平稳所需的更新次数，以及终值与种子间波动的比较。待数据确认后填写。]

\begin{figure}[t]
  \centering
  \figph{fig6_5}{五种学习型方法的训练曲线：验证集目标值随更新次数的变化，浅色为各种子，深色为种子均值的滑动平均，阴影为种子范围}
  \caption{学习型方法的收敛过程}\label{fig:converge}
\end{figure}

计算开销见表~\ref{tab:cost}，分为离线规则搜索、离线训练、在线每次决策与每实例求解四部分。规则生成是一次性的离线成本，规则库固定后在线阶段不再调用 LLM：LLM 记录所用模型与版本、调用次数、令牌数与费用；GP-Best 记录有效候选评价次数与 CPU 时间。GP 的评价预算取 LLM 实际完成的有效评价次数（6.1.3 节），两者实际完成的次数均如实列出。学习型方法另外报告模型参数量与训练期间的环境交互步数，而不只是训练回合数。滚动重排程遗传算法只在新作业到达时触发，每次搜索 3 秒，每个实例平均重排程 \tbd[※] 次；学习型方法在每启动一道工序时触发，每次决策只需一次前向计算，耗时 \tbd[※] 毫秒。两类方法的触发频率与每次决策的权力不同，遗传算法每次同时决定机器分配与全部待加工工序的顺序，规则类方法每次只决定一道工序，因此二者单次决策耗时之比不能视为完全公平的速度比较，表中以每实例求解总耗时作为主要对照。

\begin{table}[t]
  \centering
  \caption{计算开销（中心场景）}\label{tab:cost}
  \small
  \setlength{\tabcolsep}{4pt}
  \begin{tabular}{@{}lccccc@{}}
    \toprule
    方法 & 离线搜索或训练/h & 参数量 & 环境交互步数 & 每决策/ms & 每实例/s \\
    \midrule
    经典规则 & --- & --- & --- & \tbd[※] & \tbd[※] \\
    库内最优单规则 & --- & --- & --- & \tbd[※] & \tbd[※] \\
    GP-Best & \tbd[※] & --- & --- & \tbd[※] & \tbd[※] \\
    滚动重排程遗传算法 & --- & --- & --- & \tbd[※] & \tbd[※] \\
    CRG--DDQN & \tbd[※] & \tbd[※] & \tbd[※] & \tbd[※] & \tbd[※] \\
    CRG--A2C & \tbd[※] & \tbd[※] & \tbd[※] & \tbd[※] & \tbd[※] \\
    CRG--Rainbow & \tbd[※] & \tbd[※] & \tbd[※] & \tbd[※] & \tbd[※] \\
    Direct-Job PPO & \tbd[※] & \tbd[※] & \tbd[※] & \tbd[※] & \tbd[※] \\
    SC--PPO（本文） & \tbd[※] & \tbd[※] & \tbd[※] & \tbd[※] & \tbd[※] \\
    \midrule
    LLM 规则生成（3 批次） & \tbd[※] & \multicolumn{4}{l}{模型 \tbd[版本]；调用 \tbd[※] 次；令牌 \tbd[※]；费用 \tbd[※]；实际有效评价 \tbd[※] 次} \\
    GP 进化 & \tbd[※] & \multicolumn{4}{l}{实际有效评价 \tbd[※] 次；CPU 时间 \tbd[※]} \\
    \bottomrule
  \end{tabular}
\end{table}

综合上述结果：\tbd[待全部数据确认后填写小结，按 6.4 节的四条证据链说明理论特征、LLM 规则生成、在线选择和区间认证Critic各自的贡献，以及相对 GP-Best 与 Direct-Job PPO 的结果；在 9 个未见场景与两类压力测试上的稳定性；与搜索类方法在目标值和耗时上的对比。未得到统计支持的部分如实说明。]


%=====================================================================
\subsection{真实工业订单轨迹驱动的案例研究}\label{sec:industrial_case}

为检验本文方法在合成实例之外的订单结构下的可执行性与适用边界，采用公开柔性包装制造数据集 \emph{Data for Flexible Packaging Scheduling}~\cite{tarigan2025flexpackdata} 开展独立案例研究。该案例采用真实企业订单的到达日期、交期、产品规格与工艺路线；机器可选集合和工序加工时间依公开设备属性重建，累计服务承诺则由事前固定的情境协议外生设定。因而本节是\emph{真实工业订单轨迹驱动的扩展情境研究}，而非具有实测合同检查点、实测工序完成日志的端到端企业实证。真实轨迹只作最终外部测试，不能用于训练、规则筛选、调节检查点参数或反复挑选表现最好的模型。

\subsubsection{数据来源与调度实例重建}

该公开数据包含59张生产订单、15台机器、26种产品目录类型（订单实际出现11种）、7种工艺路线（实际采用6种）和6类工序；将每张订单的产品类型依次关联到工艺路线后可得到237道工序。原始订单的\texttt{InDate}位于2022年9月11日至10月1日，\texttt{Deadline}位于2022年10月15日至11月14日。生产订单具有订单编号、产品类型、优先级、加工长度（\texttt{Running Meter}）及上述日期，但不包含企业实测的逐工序开始与完工时间，也不包含真实的客户累计供货协议。为避免将源数据与研究者设定混淆，表~\ref{tab:industrial_data_mapping}逐项列出使用口径。

\begin{table}[t]
  \centering
  \caption{柔性包装公开数据与本文模型的映射及证据边界}\label{tab:industrial_data_mapping}
  \footnotesize
  \setlength{\tabcolsep}{4pt}
  \begin{tabular}{@{}p{3.1cm}p{4.1cm}p{6.4cm}@{}}
    \toprule
    模型信息 & 数据字段或来源 & 处理口径与性质 \\
    \midrule
    作业、释放时间、交期 & \texttt{Order.No}, \texttt{InDate}, \texttt{Deadline} & 真实订单字段；日期仅有天级精度，转为时间轴时固定日期边界约定 \\
    产品与工序顺序 & 产品类型、\texttt{Routing}、\texttt{Process} & 原始表连接可重建237道有序工序 \\
    机器资格 & \texttt{Machine.Code}, \texttt{Color}及工序类别 & 按预先定义的工序组和印刷色数重建；原始数据未给完整机台宽幅资格表 \\
    加工时间 $p_{j\ell m}$ & \texttt{Running Meter}, \texttt{MinSpeed}, \texttt{Ratio Speed} & 按公开速度字段估算，不当作实际加工日志；源中的设置时间不纳入当前无换线模型 \\
    服务数量 $q_j$ & \texttt{Running Meter} & 在同类可互换产品内累计；数量单位和归一化必须在测试前固定 \\
    服务实体、承诺与检查点 & 无对应原始合同字段 & 依据\emph{事前外生的模拟合同协议}构造；不能声称来自企业真实合同 \\
    \bottomrule
  \end{tabular}
\end{table}

令测试时间原点为研究协议公布的回放开始日，日期边界、时区与时间单位须在查看算法测试结果之前一致固定。订单只能在其记录的\texttt{InDate}对应时刻释放；只有被释放的订单才向调度策略暴露产品、工序、数量与交期。加工时间采用 \(p_{j\ell m}=\mathrm{RunningMeter}_j/(\mathrm{MinSpeed}_m\,\mathrm{RatioSpeed}_m)\) 这一可重建估算式；分母单位需与工厂机器表一致，机器资格按照工序组及印刷色数重建，无法从源数据确认的宽幅资格不额外虚构为真实约束。源数据的换线准备时间和设备故障未纳入本文模型，案例也不宣称重演实际企业的完工轨迹。原始表、处理脚本、字段映射及机器候选清单随复现材料公开，以便复核。

\subsubsection{外生累计承诺与独立测试协议}

为与第3节的信息边界一致，必须先于回放固定服务实体集合、合同计量单位、\(Q_r\)、\(Q_{rk}\)、\(\tau_{rk}\)、\(w_{rk}\)、公开预期份额\(\pi_r\)，以及复合泊松近似预测使用的到达强度与服务数量矩。服务实体只汇总\emph{可按同一计量口径互换履约}的产品，不能为追求更好的数值结果任意跨产品型号合并；产品目录、机器能力及回放起点已知订单属于可用于事前定义的公开信息。服务承诺可按第6.1节固定的两检查点比例与权重比例形成\emph{模拟合同}，但其总量和截止窗口应根据回放开始时已公开的合同情境或独立的外部先验确定；若缺乏外部历史数据，应在正式测试前登记一套固定参数，并明确标记为研究者设定。所有候选合同情境及其物理可行性应先于算法比较确定并公布，不得根据本次59张订单中\emph{尚未释放}的总加工长度、未来产品出现次数、未来截止日期或SC--PPO测试结果反向设定。

特别地，\(Q_r=\sum_{j:e_j=r}q_j\)（对整条测试轨迹求和）并非合格的初始已知合同构造方式，因为它会泄漏尚未释放订单的具体总供给；由整条轨迹实际频数估计\(\pi_r\)或到达率再在\(t=0\)提供给Actor也存在相同问题。真实订单实际到达不需要服从泊松过程，仍沿用第4节的\emph{复合泊松在线近似预测}并独立报告预测失配；不允许以完整测试轨迹拟合后把未来参数作为已知输入。若日期级释放与原合成环境、米数级\(q_j\)或15台机器的输入维度不匹配，采用固定、可复现的单位转换及变规模编码适配，不得借测试集成绩调节归一化常数；发生迁移训练或微调时另作报告，不归类为直接零样本测试。

按照两个层次评价：(i) 利用真实订单和重建加工资源检查全部方法的调度合法性及TT、按时完工率与运行时延；(ii) 在同一真实轨迹上叠加事前公布的模拟累计合同，比较TT、WSF、综合目标\(F\)、检查点满足率与规则切换行为。前者属于真实轨迹下的作业层证据，后者属于合同扩展条件下的服务层证据，两种证据不得混称为真实企业合同履约效果。案例沿用已经确定的训练末期模型与规则库，以EDD、ATC、Best-Single、Direct-Job PPO及本文方法作配对比较；对CC--DVL候选模型，同时报告原标量证书Critic（C4），避免将新价值机制的潜在收益与规则动作抽象混为一谈。

\begin{table}[t]
  \centering
  \caption{真实订单轨迹与模拟累计合同下的案例对比（结果待独立实验填入）}\label{tab:industrial_case_result}
  \footnotesize
  \setlength{\tabcolsep}{4pt}
  \begin{tabular}{@{}lcccccc@{}}
    \toprule
    方法 & TT & WSF & $F$ & 首检查点满足率 & 最终满足率 & 决策耗时/ms \\
    \midrule
    EDD & \tbd[待测] & \tbd[待测] & \tbd[待测] & \tbd[待测] & \tbd[待测] & \tbd[待测] \\
    ATC & \tbd[待测] & \tbd[待测] & \tbd[待测] & \tbd[待测] & \tbd[待测] & \tbd[待测] \\
    Best-Single & \tbd[待测] & \tbd[待测] & \tbd[待测] & \tbd[待测] & \tbd[待测] & \tbd[待测] \\
    Direct-Job PPO & \tbd[待测] & \tbd[待测] & \tbd[待测] & \tbd[待测] & \tbd[待测] & \tbd[待测] \\
    SC--PPO（C4） & \tbd[待测] & \tbd[待测] & \tbd[待测] & \tbd[待测] & \tbd[待测] & \tbd[待测] \\
    SC--PPO + CC--DVL & \tbd[待测] & \tbd[待测] & \tbd[待测] & \tbd[待测] & \tbd[待测] & \tbd[待测] \\
    \bottomrule
  \end{tabular}
\end{table}

\subsubsection{履约轨迹与管理解释}

对一个\emph{事先按规则选定}的服务实体，绘制不同方法的累计完成量\(Y_{rt}\)与公开检查点目标\(Q_{rk}\)；同图标注真实订单释放事件以及各检查点结算时刻，以分析是否因提前调配有限产能而减少阶段短缺（图~\ref{fig:industrial_service_trace}）。优先按事前约定的合同规模或检查点紧迫度选取展示对象，不以本方法改进幅度反向挑选；同时报告所有实体的总体指标，避免单条优势轨迹代表总体结论。可以在同一真实到单轨迹上使用预注册的宽松/中性/紧张模拟合同进行敏感性比较，但这些情境共享同一真实订单样本，不得当作多条独立工业轨迹用于统计显著性检验。

\begin{figure}[t]
  \centering
  \figph{fig6_industrial_trace}{真实柔性包装订单的累计完成数量与模拟合同检查点：横轴为物理时间，纵轴为同产品履约数量，叠加EDD、Best-Single、C4及CC--DVL的实际回放轨迹；注明原始释放事件与检查点时刻（待完成算法测试后绘制）}
  \caption{真实工业订单轨迹驱动案例的累计服务履约图（待实验生成）}\label{fig:industrial_service_trace}
\end{figure}

本案例仅包含一条真实企业订单轨迹：不同训练种子和不同合同情境并非相互独立的企业样本。结论应限于所公开生产订单、可审计的工艺与加工时间重建假设，以及明确列出的模拟合同条件；不能由单个案例推断跨企业的普遍改进。\tbd[待实验填写：相同订单轨迹下各算法TT--WSF权衡、检查点履约情况、零样本迁移是否可行与计算耗时；若CC--DVL未优于C4，应如实说明。]